% STATUS: written 2026-10-12
% Sources read in full: Herstein 1978 (pp. 369--370), Herstein 1979 (pp. 509--511), Bresar 1993 (pp. 385--388, introduction and Section 2).
\section{The school of Herstein and the commutativity theorems}\label{sec:prime-semiprime:herstein-school}

Each theorem below is a statement about derivations of a prime or semiprime ring, which is what the school of Herstein proved.

Section~\ref{sec:prime-semiprime:posner} ended with two theorems whose
conclusion is that a prime ring is commutative. We ask which other conditions
on the values of a derivation give the same conclusion. Herstein took up this
question in two short notes \cite{Herstein1978,Herstein1979}.

The first note proves two results about derivations \cite[p.~369]{Herstein1978}. Herstein could not find them in the literature. We state the first result, which holds in any ring (Theorem~\ref{thm:prime:herstein-ideal}).

\begin{theorem}[{\cite[Theorem~1, p.~369]{Herstein1978}}]\label{thm:prime:herstein-ideal}
Let $R$ be a ring and let $d$ be a derivation of $R$ with $d^3\neq 0$. Then the
subring generated by all values $d(r)$, for $r\in R$, contains a nonzero ideal
of $R$.
\end{theorem}

The condition $d^3\neq 0$ cannot be dropped, and Herstein gives an example with
an inner derivation.

\begin{example}[{\cite[p.~369]{Herstein1978}}]\label{ex:prime:herstein-cube}
Let $R$ be a prime ring with a nonzero element $a$ such that $a^2=0$. Define
\[
d(x)=ax-xa .
\]
The subring $B=aR+Ra$ contains $d(R)$, and $d^3=0$. Yet $B$ contains no nonzero
ideal of $R$, because $aBa=0$.
\end{example}

We state the second result, a commutativity theorem for prime rings.

\begin{theorem}[{\cite[Theorem~2, p.~369]{Herstein1978}}]\label{thm:prime:herstein-commuting}
Let $R$ be a prime ring and let $d\neq 0$ be a derivation of $R$ such that
$d(x)d(y)=d(y)d(x)$ for all $x,y\in R$. If $\operatorname{char}R\neq 2$, then
$R$ is a commutative integral domain. If $\operatorname{char}R=2$, then $R$ is
commutative or an order in a simple algebra of dimension four over its center.
\end{theorem}

Here the values of $d$ commute with each other. In
Theorem~\ref{thm:prime:posner-two} the condition links $d(a)$ with $a$ itself.
Herstein asks whether Theorem~\ref{thm:prime:herstein-commuting} extends from
commuting values to values that satisfy the standard identity of degree $k$
\cite[p.~370]{Herstein1978}. He leaves this as a question.

The second note turns the question around \cite[p.~509]{Herstein1979}. It asks which elements of a prime ring commute with all the values of a nonzero derivation, and Herstein settles it (Theorem~\ref{thm:prime:herstein-commute-values}) \cite[p.~509]{Herstein1979}.

\begin{theorem}[{\cite[Theorem, p.~509]{Herstein1979}}]\label{thm:prime:herstein-commute-values}
Let $R$ be a prime ring and let $d\neq 0$ be a derivation of $R$. Suppose that
$a\in R$ is such that $ad(x)=d(x)a$ for all $x\in R$. If $R$ is not of
characteristic $2$, then $a$ is in the center $Z$ of $R$. If $R$ has
characteristic $2$, then $a^2\in Z$, and if $a\notin Z$ then $d(x)=(\lambda
a)x-x(\lambda a)$ for all $x\in R$, where $\lambda$ lies in the extended centroid
of $R$.
\end{theorem}

Away from characteristic $2$ the answer is simple: $a$ is central
\cite[p.~509]{Herstein1979}. In characteristic $2$ the derivation is inner, and
the element that induces it lies in the extended centroid
\cite[p.~509]{Herstein1979}. The extended centroid also appeared in the
characteristic $2$ case of Section~\ref{sec:prime-semiprime:posner}.
Section~\ref{sec:prime-semiprime:quotient-rings} explains where it lives.

Bre\v{s}ar recalls the theorem of Posner: a prime ring with a nonzero
centralizing derivation is commutative \cite[p.~385]{Bresar1993}. He adds that
Mayne proved the analogous result for centralizing automorphisms
\cite[p.~113]{Mayne1976Centralizing}. A map $F$ of a
ring is \emph{centralizing} on a set $S$ of elements if $[s,F(s)]$ is central
for all $s\in S$, and \emph{commuting} if $[s,F(s)]=0$.

Bre\v{s}ar lists many authors who extended both theorems to further maps and
sets of elements \cite[p.~385]{Bresar1993}. He then describes all additive maps
that are centralizing on a prime ring (Theorem~\ref{thm:prime:bresar-centralizing}).

\begin{theorem}[{\cite[Theorem~A, p.~385]{Bresar1993}}]\label{thm:prime:bresar-centralizing}
Let $R$ be a prime ring and let $F$ be an additive map of $R$ into itself that
is centralizing on $R$. If $R$ has characteristic different from $2$, or if $F$
is commuting on $R$, then $F(x)=\lambda x+\zeta(x)$ for $x\in R$, where $\lambda$
is an element of the extended centroid $C$ of $R$ and $\zeta$ is an additive map
of $R$ into $C$.
\end{theorem}

Theorem~\ref{thm:prime:posner-two} is the case of a derivation. Bre\v{s}ar also
treats two derivations at once (Theorem~\ref{thm:prime:bresar-two}) \cite[Theorem~B, p.~386]{Bresar1993}.

\begin{theorem}[{\cite[Theorem~B, p.~386]{Bresar1993}}]\label{thm:prime:bresar-two}
Let $R$ be a prime ring and let $U$ be a nonzero left ideal of $R$. Let $d$ and
$g$ be derivations of $R$ such that $d(u)u-ug(u)$ is central for all $u\in U$.
If $d\neq 0$, then $R$ is commutative.
\end{theorem}

Bell and Martindale had proved the case of one derivation when $U$ is a
one-sided ideal \cite[p.~386]{Bresar1993}.

The extended centroid enters the proofs of Bre\v{s}ar when an identity such as
$axb=bxa$ for all $x$ forces $b=\lambda a$ \cite[Remark~6, p.~387]{Bresar1993}.
Section~\ref{sec:prime-semiprime:quotient-rings} turns to the ring in which such
a $\lambda$ lives.

A generalized derivation whose values are nilpotent, power-central or invertible is shown to force the ring to commute, in the line of the Herstein school of the chapter~\cite{WlGeneralizedDerivationsWith2019}.

A generalized derivation whose values are nilpotent, power-central or invertible is shown to force the ring to commute, in the line of the Herstein school~\cite{WlGeneralizedDerivationsWithNilp2019}.
